\begin{frame}{5.1 Statistik Gerbang dan Sekuensial I}
    Sirkuit $4 \times 4$ memerlukan total \textbf{28 operasi gerbang} sekuensial.
    \textbf{Sekuensial I: Pre-kondisi Register Target ($q_1, q_2$)}
    \begin{itemize}
        \item \textbf{Operasi Fase $P(\pi/4)$ pada $q_1$:} Menyiapkan korelasi internal.
        \item \textbf{Gerbang Terkontrol $CU$ (1b):} Rotasi uniter antar target dengan $\theta=1.1747, \phi=-2.83038, \lambda=3.83087$:
        \begin{equation*}
            U_{1b} \approx \begin{pmatrix} 0.832 & 0.425 + i0.357 \\ -0.528 + i0.170 & 0.449 + i0.700 \end{pmatrix}
        \end{equation*}
    \end{itemize}
\end{frame}

\begin{frame}{5.2 Sekuensial II dan Refinement Fase}
    \textbf{Sekuensial II: Inisiasi Coupling ke Register Fase ($q_0$)}
    \begin{itemize}
        \item \textbf{Gerbang $CX(q_0, q_1)$:} Menciptakan \textit{entanglement} antara fase dan target.
        \item \textbf{Koreksi Fase $P(-\pi/4)$:} Mengimbangi rotasi sebelumnya.
        \item \textbf{Gerbang $CU$ (2c):} $\theta=1.1747, \phi=-0.689273, \lambda=-0.31121$.
    \end{itemize}
    \textbf{Refinement (20 Operasi Sisa):}
    Terdiri dari 10x $CU$, 5x $P$, dan 4x $CX$ untuk menyelaraskan fasa maturitas 1, 3, dan 6 bulan serta memitigasi fasa parasit (\textit{stray phases}).
\end{frame}

\begin{frame}{5.3 IQFT 1-bit dan Mekanisme Interferensi}
    Transformasi balik dilakukan dengan gerbang Hadamard pada $q_0$:
    \begin{equation}
        |\Phi_{final}\rangle = \frac{1}{2} \left[ (1 + e^{2\pi i \phi})|0\rangle + (1 - e^{2\pi i \phi})|1\rangle \right] \otimes |\psi_4\rangle
    \end{equation}
    \textbf{Interpretasi Pengukuran pada $q_0$:}
    \begin{itemize}
        \item \textbf{Fase $\phi \approx 0$:} Amplitudo state $|0\rangle$ berinterferensi konstruktif ($1+1=2$).
        \item \textbf{Fase $\phi \approx 0.5$:} Amplitudo state $|1\rangle$ menjadi konstruktif ($1 - (-1) = 2$), mengindikasikan \textit{eigenvalue} yang signifikan.
    \end{itemize}
    Hasil ini divisualisasikan melalui histogram QASM yang menunjukkan indikator kasar besaran \textit{eigenvalue} utama.
\end{frame}
